\documentclass[a4paper, 12pt]{exam}
\usepackage[UTF8]{ctex}
\usepackage{amsmath}
\usepackage{amssymb}
\usepackage{graphicx}
\usepackage{xeCJK}
\usepackage[a4paper, margin=2cm]{geometry}
\renewcommand{\thepartno}{\arabic{partno}}
\pointpoints{分}{分}
\setlength{\parindent}{2em}  % 每段缩进2个字宽
\renewcommand{\baselinestretch}{1.5}

\begin{document}
\title{植物生物学实验试题}
\author{苗老师/杨老师}
\maketitle
% \chapter*{实验：植物学}
\newpage
\noindent{\large{\textbf{实验：植物学（40分）}}}
\begin{questions}
    \question[10]{写出示范台上各显微镜指针所指的构造名称}
    \begin{parts}
        \part{\rule{3cm}{0.2pt}}
        \part{\rule{3cm}{0.2pt}}
        \part{\rule{3cm}{0.2pt}}
        \part{\rule{3cm}{0.2pt}}
    \end{parts}
    \question[30]{实验材料为寄生植物菟丝子，完成下列实验内容}
    \begin{parts}
        \part[10]{解剖菟丝子的花，写出花公式，绘其花图式，判断属哪个科和传粉类型，说明理由。}
        \vfil
        \part[10]{分别做菟丝子的茎和寄主（茎或叶）的横切，绘其结构简图，并注示各部组成的名称。}
        \vfil
        \newpage
        \part[3]{苑丝子是通过什么结构附着并侵入宿主体内的？并绘简图表示。}
        \vfil
        \part[4]{分析菟丝子与寄主的关系，判断以下叙述是否正确，简要说明理由}
        \begin{subparts}
            \subpart{菟丝子抑制寄主光合作用；}
            \subpart{菟丝子仅通过土壤与寄主竞争水分。}
            \vfil
        \end{subparts}
        \part[3]{在生产实践中，菟丝子被用于防治入侵植物加拿大一枝黄花。试分析菟丝子导致寄主生物量变化的机制。}
        \vfil
        \newpage
    \end{parts}

\end{questions}

{\large{\textbf{实验：植物生理学（60分）}}}

{\large{\textbf{一、水势测定（50分）}}}

水势(Waterpotential)能够表征植物组织吸收水的能力。水势的值取决于植物的细胞质基质与蒸馏水之间的势差。水会从水势高的区域流向水势低的区域。当区域间的水势差越大，水会更快地从水势高的区域流向水势低的区域。也就是说，当水势梯度越大时，水的流失度会越快。

组织中的渗透势可以通过范特霍夫方程(Van't Hoff equation)计算得出。范特霍夫方程如下:
\[\Psi=-\mathrm{iCRT}\]

其中，R表示通用气体常数，即8.31J/mol*K;T表示温度，即273℃+t，假定本实验在常温下完成，即t=22℃；i表示 Van't Hoff等渗系数，对于所有非电解质，例如蔗糖，i=1；C表示等渗溶液的浓度。

实验必要物品包括：（1）土豆；（2）刀片；（3）浓度为1M的蔗糖溶液，体积为50mL；（4）蒸馏水；（5）移液器及吸耳球；（6）6支15mL试管；（7）圆头镊子；（8）吸水纸；（9）尺子；（10）天平。
\vspace{2em}
本实验步骤如下：
\begin{questions}
    \question[6]{请用浓度为1M的蔗糖溶液，分别在6支试管中配制浓度为不同浓度梯度的蔗糖溶液 \textbf{（请自行设计）}，每支试管中蔗糖溶液的体积均为10mL。请在答卷上注明，配制上述蔗糖溶液需要浓度为1M的蔗糖溶液的体积是多少?需要蒸馏水的体积是多少?}
    \begin{table}[htbp]
        \centering
        \begin{tabular}{
                |c|c|c|c|c|c|c|c|c|}
            \hline
            溶液（M）      &    &    &    &    &    &    & 总所需（mL） \\
            \hline

            1M Suc（mL） &    &    &    &    &    &    &         \\
            \hline

            蒸馏水（mL）    &    &    &    &    &    &    &         \\
            \hline

            总体积（mL）    & 10 & 10 & 10 & 10 & 10 & 10 & 60      \\
            \hline
        \end{tabular}
    \end{table}
    \question[6]{请用解剖刀片将土豆切为6块，每块土豆的长度约为20至30mm，宽度约为5mm，高度约为5mm（即\(5\rm{mm}\times5\rm{mm}\)的横截面）。请注意，土豆切块时间尽量短，以避免土豆块失水变干。请测量每块土豆的具体长度（精确到小数点后1位）与重量（精确到小数点后2位），并将数据记录。}
    \begin{table}[htbp]
        \centering
        \begin{tabular}{|p{5em}|p{5em}|p{5em}|p{5em}|p{5em}|p{5em}|}
            \hline
            溶液（M）  &  &  &  &  & \\
            \hline
            长度（mm） &  &  &  &  & \\
            \hline
            重量（g）  &  &  &  &  & \\
            \hline
        \end{tabular}
    \end{table}
    \newpage
    \question{请将切好的土豆块放置于上述不同的蔗糖溶液的试管中（请分别将每个主豆块置于1支试管中），静置20分钟。}
    \question[6]{静置20分钟后，请取出土豆块，请重新测量每块土豆的具体长度（精确到小数点后1位）与重量（精确到小数点后2位），并将数据记录。}
    \begin{table}[htbp]
        \centering
        \begin{tabular}{|p{5em}|p{5em}|p{5em}|p{5em}|p{5em}|p{5em}|}
            \hline
            溶液（M）  &  &  &  &  & \\
            \hline
            长度（mm） &  &  &  &  & \\
            \hline
            重量（g）  &  &  &  &  & \\
            \hline
        \end{tabular}
    \end{table}
    \question[3]{请利用你测量的结果，得出等渗溶液的浓度C，单位为M。请将结果记录。

        等渗浓度为：}
    \question[8]{请利用范特霍夫方程\(\Psi=-\rm{iCRT}\)计算土豆块的水势。请将计算结果记录。计算过程与计算结果要求写出单位。最终结果要求写出压强单位以及国际单位制的单位。

        *国际单位制的单位包括米(m)，千克(kg)，秒(S)，安培(A)，开尔文(K)摩尔(mol)，坎德拉(cd)。}
\end{questions}
\end{document}